\documentclass[aspectratio=169]{beamer}
\usetheme{metropolis}
\usepackage{amsmath, amssymb}
\usepackage{tikz}
\usetikzlibrary{positioning, arrows.meta}
\usepackage{pgfplots}
\pgfplotsset{compat=1.18}
\usepackage{luatexko}
\setmainhangulfont{NanumGothic}
\setsanshangulfont{NanumGothic}
\usepackage{xcolor}
\usepackage{booktabs}

% 색상 정의
\definecolor{primary}{RGB}{20, 40, 80}
\definecolor{accent}{RGB}{220, 80, 60}
\definecolor{lightgray}{RGB}{240, 240, 240}
\definecolor{darkgray}{RGB}{60, 60, 60}

\setbeamercolor{frametitle}{bg=primary, fg=white}
\setbeamercolor{title separator}{fg=accent}
\setbeamercolor{progress bar}{fg=accent}
\setbeamercolor{alerted text}{fg=accent}

\title{Universality in Phase Transition}
\subtitle{Lattice Gas vs. Molecular Dynamics}
\author{이호준 (25-095)}
\institute{탐구물리 개인 프로젝트}
\date{June 2, 2026}

\begin{document}

% ── 슬라이드 1: 제목 ──────────────────────────────────────────
\maketitle

% ── 슬라이드 2: 동기 ──────────────────────────────────────────
\begin{frame}{동기: 같은 현상, 다른 해상도}

\begin{columns}[T]
\begin{column}{0.48\textwidth}
  \centering
  \textbf{\large Microscopic (MD)}\\[0.4em]
  \begin{tikzpicture}[scale=0.85]
    % 상자
    \draw[thick, primary] (0,0) rectangle (3.5,3.5);
    % 입자들 (뭉친 상태 - 액체)
    \foreach \x/\y in {
      0.6/0.7, 1.0/0.5, 0.8/1.1, 1.3/0.8, 1.1/1.4,
      0.5/1.5, 1.5/1.1, 0.7/0.3, 1.2/0.2, 1.6/0.5,
      0.9/1.8, 1.4/1.7, 0.6/2.0, 1.0/2.1}{
      \fill[accent!80] (\x,\y) circle (0.13);
    }
    % 기체 입자들
    \foreach \x/\y in {2.5/0.4, 3.0/1.2, 2.8/2.5, 3.2/3.0, 2.2/3.1}{
      \fill[primary!50] (\x,\y) circle (0.10);
    }
    \node[below, darkgray] at (1.75, -0.1) {\small 입자 $N$개, $F=ma$};
  \end{tikzpicture}\\[0.3em]
  {\small Lennard-Jones potential}\\
  {\small 입자 하나하나를 추적}
\end{column}

\begin{column}{0.04\textwidth}
  \centering
  \vspace{2cm}
  \Large$\Leftrightarrow$
\end{column}

\begin{column}{0.48\textwidth}
  \centering
  \textbf{\large Macroscopic (Lattice Gas)}\\[0.4em]
  \begin{tikzpicture}[scale=0.85]
    \draw[thick, primary] (0,0) rectangle (3.5,3.5);
    % 격자
    \foreach \i in {0,...,6}{
      \draw[lightgray!80] (0, \i*0.5) -- (3.5, \i*0.5);
      \draw[lightgray!80] (\i*0.5, 0) -- (\i*0.5, 3.5);
    }
    % 뭉친 입자 (점유 칸)
    \foreach \i/\j in {
      1/1, 2/1, 1/2, 2/2, 3/1, 1/3, 2/3, 3/2,
      0/1, 0/2, 1/0, 2/0}{
      \fill[accent!70] (\i*0.5+0.05, \j*0.5+0.05)
        rectangle (\i*0.5+0.45, \j*0.5+0.45);
    }
    % 빈 칸 (기체)
    \foreach \i/\j in {4/4, 5/3, 6/5, 4/6, 6/2}{
      \fill[primary!20] (\i*0.5+0.05, \j*0.5+0.05)
        rectangle (\i*0.5+0.45, \j*0.5+0.45);
    }
    \node[below, darkgray] at (1.75, -0.1) {\small 격자 $N\times N$};
  \end{tikzpicture}\\[0.3em]
  {\small Metropolis algorithm}\\
  {\small 칸별 점유 여부만 추적}
\end{column}
\end{columns}

\vspace{0.5em}
\centering
\alert{\textbf{Research Question:} 두 모델은 같은 phase transition을 예측하는가?}

\end{frame}

% ── 슬라이드 3: 이론 ──────────────────────────────────────────
\begin{frame}{이론 배경}

\begin{columns}[T]
\begin{column}{0.48\textwidth}

\textbf{Lennard-Jones Potential (MD)}
\begin{equation*}
V(r) = 4\varepsilon\left[\left(\frac{\sigma}{r}\right)^{12} - \left(\frac{\sigma}{r}\right)^{6}\right]
\end{equation*}

\begin{tikzpicture}[scale=0.9]
\begin{axis}[
  width=6cm, height=4cm,
  xmin=0.8, xmax=3.0, ymin=-1.3, ymax=2.0,
  xlabel={$r/\sigma$}, ylabel={$V/\varepsilon$},
  axis lines=left,
  tick label style={font=\tiny},
  label style={font=\small},
  xtick={1,1.5,2,2.5,3},
  ytick={-1,0,1,2},
]
\addplot[accent, thick, domain=0.9:3.0, samples=100]
  {4*(x^(-12) - x^(-6))};
\draw[dashed, darkgray] (axis cs:1.122,-1.0) -- (axis cs:1.122,0.2);
\node[font=\tiny] at (axis cs:1.5,-1.1) {$r_{min}$};
\draw[<->, darkgray, thin] (axis cs:0.95,-1.0) -- (axis cs:0.95,0);
\node[font=\tiny, left] at (axis cs:0.92,-0.5) {$\varepsilon$};
\end{axis}
\end{tikzpicture}

\vspace{0.3em}
{\small 인력 $\varepsilon$ vs 열에너지 $k_BT$ → 상전이}

\end{column}
\begin{column}{0.48\textwidth}

\textbf{Lattice Gas = Ising Model}
\begin{equation*}
E = -J\sum_{\langle i,j\rangle} s_i s_j, \quad s_i \in \{-1,+1\}
\end{equation*}

\vspace{0.3em}
\textbf{Onsager의 정확한 해석해 (1944):}
\begin{equation*}
\alert{T_c = \frac{2J}{k_B\ln(1+\sqrt{2})} \approx 2.269\,\frac{J}{k_B}}
\end{equation*}

\vspace{0.5em}
\textbf{Universality (보편성):}\\[0.3em]
임계점 근처에서 임계지수 $\beta$는\\
미시적 디테일에 \alert{무관}
\begin{equation*}
M \propto (T_c - T)^{\beta}, \quad \beta = \tfrac{1}{8}
\end{equation*}

\end{column}
\end{columns}

\end{frame}

% ── 슬라이드 4: 방법론 1 - 격자기체 ─────────────────────────
\begin{frame}{방법론 1: Lattice Gas (Metropolis MC)}

\begin{columns}[T]
\begin{column}{0.38\textwidth}

\textbf{Metropolis Algorithm}\\[0.4em]

\begin{tikzpicture}[
  node distance=1.0cm,
  box/.style={draw=primary, rounded corners=3pt,
    fill=primary!10, font=\scriptsize,
    text width=3.0cm, align=center, minimum height=0.6cm},
  decision/.style={draw=accent, rounded corners=3pt,
    fill=accent!10, font=\scriptsize,
    text width=3.0cm, align=center, minimum height=0.6cm},
  arrow/.style={->, thick, darkgray}
]
  \node[box] (init) {격자 초기화 (random)};
  \node[box, below of=init] (pick) {스핀 $(i,j)$ 랜덤 선택};
  \node[box, below of=pick] (de) {$\Delta E = 2Js_{ij}\!\sum_\text{nbr}\! s$};
  \node[decision, below of=de] (check) {$\Delta E \leq 0$?};
  \node[box, below of=check] (yes) {Yes: 뒤집기};
  \node[box, below of=yes] (no) {No: 확률 $e^{-\Delta E/k_BT}$로};
  \node[box, below of=no] (meas) {$M$, $E$ 측정};

  \draw[arrow] (init) -- (pick);
  \draw[arrow] (pick) -- (de);
  \draw[arrow] (de) -- (check);
  \draw[arrow] (check) -- node[right, font=\tiny]{Yes} (yes);
  \draw[arrow] (check.east) -- ++(0.4,0)
    |- node[right, font=\tiny, pos=0.25]{No} (no.east);
  \draw[arrow] (yes) -- (no);
  \draw[arrow] (no) -- (meas);
  \draw[arrow] (meas.west) -- ++(-0.3,0)
    |- node[left, font=\tiny]{반복} (pick.west);
\end{tikzpicture}

\end{column}
\begin{column}{0.59\textwidth}

\textbf{측정량 \& $T_c$ 찾기}\\[0.4em]

\begin{tikzpicture}
\begin{axis}[
  width=7.0cm, height=3.0cm,
  xmin=1.5, xmax=3.5,
  xlabel={$T$}, ylabel={$|M|$},
  axis lines=left,
  tick label style={font=\tiny},
  label style={font=\small},
  xtick={1.5,2.269,3.5},
  xticklabels={, $T_c$, },
  ytick={0,0.5,1},
]
\addplot[accent, thick, domain=1.5:2.269, samples=50]
  {(1 - (sinh(2*2.269/x))^(-4))^(0.125)};
\addplot[accent, thick, domain=2.269:3.5, samples=50] {0};
\draw[dashed, darkgray] (axis cs:2.269,0) -- (axis cs:2.269,1);
\end{axis}
\end{tikzpicture}

\vspace{0.3em}
\begin{tabular}{p{1.5cm}p{3.2cm}p{2.5cm}}
\toprule
측정량 & 계산 & 목적\\
\midrule
$|M|(T)$ & 스핀 평균 & $T_c$ 위치\\[0.2em]
$C_v(T)$ & $\frac{\langle E^2\rangle-\langle E\rangle^2}{k_BT^2}$ & $T_c$ 봉우리\\[0.4em]
$U_L(T)$ & Binder cumulant & $T_c$ 정밀 교점\\
\bottomrule
\end{tabular}

\vspace{0.35em}
{\small \textbf{임계지수 추출:}
$\log|M|$ vs $\log(T_c-T)$ 기울기 $\to \beta$}

\end{column}
\end{columns}

\end{frame}

% ── 슬라이드 5: 방법론 2 - MD ────────────────────────────────
\begin{frame}{방법론 2: Molecular Dynamics (LJ Fluid)}

\begin{columns}[T]
\begin{column}{0.48\textwidth}

\textbf{시뮬레이션 설정}\\[0.4em]

\begin{itemize}
  \item $N$개 LJ 입자, 주기 경계 조건
  \item Velocity Verlet 시간 적분
  \item Thermostat으로 온도 $T$ 제어
  \item 환산단위: $\varepsilon$, $\sigma$, $m = 1$
\end{itemize}

\vspace{0.5em}
\textbf{Coexistence Curve}\\[0.3em]

\begin{tikzpicture}[scale=0.95]
\begin{axis}[
  width=6cm, height=4cm,
  xmin=0, xmax=1.0,
  ymin=0.8, ymax=1.45,
  xlabel={밀도 $\rho^*$},
  ylabel={온도 $T^*$},
  axis lines=left,
  tick label style={font=\tiny},
  label style={font=\small},
]
\addplot[accent, thick, domain=0.05:0.85, samples=60]
  {1.326 - 1.2*(x-0.32)^2 - 2.5*(x-0.32)^4};
\addplot[mark=*, mark size=3pt, primary] coordinates {(0.32, 1.326)};
\node[font=\tiny, above right] at (axis cs:0.32, 1.326) {Critical point};
\node[font=\tiny, accent] at (axis cs:0.10, 0.92) {gas};
\node[font=\tiny, accent] at (axis cs:0.72, 0.92) {liquid};
\end{axis}
\end{tikzpicture}

\end{column}
\begin{column}{0.48\textwidth}

\textbf{$T_c$ \& $\beta$ 추출 방법}\\[0.5em]

온도별로 $\rho_l$, $\rho_g$ 측정\\[0.3em]

\begin{equation*}
\rho_l - \rho_g \propto (T_c - T)^{\beta}
\end{equation*}

\vspace{0.3em}

\begin{tikzpicture}[scale=0.9]
\begin{axis}[
  width=6cm, height=3.5cm,
  xmin=-2.5, xmax=0,
  ymin=-2.5, ymax=0,
  xlabel={$\log(T_c - T)$},
  ylabel={$\log(\rho_l - \rho_g)$},
  axis lines=left,
  tick label style={font=\tiny},
  label style={font=\small},
]
\addplot[accent, thick, domain=-2.3:-0.1, samples=30]
  {0.125*x - 0.5};
\node[font=\small, primary] at (axis cs:-1.8,-1.5)
  {기울기 $= \beta$};
\end{axis}
\end{tikzpicture}

\vspace{0.3em}
{\small 로그-로그 그래프 기울기에서 $\beta$ 추출}\\[0.3em]
{\small 환산온도 $t = (T-T_c)/T_c$ 로 무차원화 후 비교}

\end{column}
\end{columns}

\end{frame}

% ── 슬라이드 6: 기대 결과 ────────────────────────────────────
\begin{frame}{기대 결과: Universality}

\begin{columns}[T]
\begin{column}{0.55\textwidth}

\begin{tikzpicture}
\begin{axis}[
  width=7.5cm, height=5.0cm,
  xmin=-0.5, xmax=0,
  ymin=-3.5, ymax=0.5,
  xlabel={$\log|t|$ \enspace ($t = (T-T_c)/T_c$)},
  ylabel={$\log(\text{order param.})$},
  axis lines=left,
  tick label style={font=\small},
  label style={font=\small},
  legend columns=1,
  legend style={
    font=\tiny,
    fill=white,
    draw=darkgray,
    at={(0.97,0.03)},
    anchor=south east,
  },
]
\addplot[accent, thick, domain=-0.45:-0.02, samples=30]
  {0.125*x - 0.3};
\addlegendentry{Lattice Gas}

\addplot[primary, thick, dashed, domain=-0.45:-0.02, samples=30]
  {0.125*x - 1.5};
\addlegendentry{MD}

\addplot[darkgray, dotted, thick, domain=-0.45:-0.02, samples=30]
  {0.125*x - 0.9};
\addlegendentry{Onsager ($\beta\!=\!1/8$)}
\end{axis}
\end{tikzpicture}

\end{column}
\begin{column}{0.42\textwidth}

\vspace{1em}
\textbf{검증 구조}\\[0.6em]

\begin{tabular}{ll}
\toprule
& $\beta$ \\
\midrule
Lattice Gas & $\approx 1/8$ \\
MD & $\approx ?$ \\
Onsager (exact) & $= 1/8$ \\
\bottomrule
\end{tabular}

\vspace{1em}
\alert{두 모델의 $\beta$가 일치한다면}\\[0.3em]
→ Universality 확인\\[0.3em]
→ 미시적 디테일과 무관하게\\
\quad 임계 거동이 결정됨을 증명

\vspace{0.5em}
{\small \textbf{단위 문제 해결:} 환산온도 $t$로 무차원화하여 같은 축에서 비교}

\end{column}
\end{columns}

\end{frame}

% ── 슬라이드 7: 방법론 요약 ──────────────────────────────────
\begin{frame}{방법론 요약}

\begin{columns}[T]
\begin{column}{0.48\textwidth}

\textbf{Lattice Gas (Metropolis MC)}\\[0.5em]
\begin{enumerate}\itemsep0.3em
  \item $N\times N$ 격자, 스핀 $s_i \in \{-1,+1\}$ 초기화
  \item Metropolis 알고리즘으로 평형 도달
  \item $|M|(T)$, $C_v(T)$, Binder cumulant $U_L$ 측정
  \item $U_L$ 교점으로 $T_c$ 확정
  \item 로그-로그 피팅으로 $\beta$ 추출\\
        → Onsager $T_c = 2.269$, $\beta = 1/8$ 과 비교
\end{enumerate}

\end{column}
\begin{column}{0.48\textwidth}

\textbf{Molecular Dynamics (LJ Fluid)}\\[0.5em]
\begin{enumerate}\itemsep0.3em
  \item $N$개 LJ 입자, Velocity Verlet 적분
  \item Thermostat으로 온도 $T^*$ 스캔
  \item 온도별 $\rho_l$, $\rho_g$ 측정 (coexistence curve)
  \item $\rho_l - \rho_g \propto (T_c - T)^\beta$ 피팅으로\\
        $T_c^*$, $\beta$ 추출
\end{enumerate}

\vspace{0.8em}
\textbf{비교}\\[0.4em]
환산온도 $t = (T-T_c)/T_c$ 로 무차원화\\[0.3em]
\alert{$\beta_\text{LG}$ vs $\beta_\text{MD}$ 일치 여부 → Universality}

\end{column}
\end{columns}

\end{frame}

\end{document}
